\section{Introduction}
\label{sec:introduction}

\subsection{To throw a stone}
To throw a stone is to make a reliable difference to the world: after the throw, the world is measurably different from how it would otherwise have been. The sentence sounds trivial until one asks what must already be in place for it to mean anything. In a substrate of micro-interactions, where is the throw? Which degrees of freedom are the thrower, which the stone, and which the outside world? When does a subsystem have \emph{choices}, rather than simply being pushed around?

This paper treats agency as something \emph{induced} rather than assumed. We do not start from goals, intentions, or preferences. We start from a finite stochastic system and ask two questions: when does a stable, budgeted subsystem exist at all, and when do its feasible choices change the distribution of outside futures? The stone is a reminder that the second question is about counterfactual difference-making under constraints: there should be at least two affordable ways of acting that lead to distinguishable outcomes.

\subsection{From objects to agents in Six Birds Theory}
In \SBT\ \citep{six_birds_theory}, macroscopic objects are not primitives. The theory proposes six emergence mechanisms, the ``six birds'', that together make stable macroscopic descriptions possible: operator rewriting, constraints on what is feasible, protocol cycles, identity and staging, closure, and resource accounting. On this view an object is a stabilised package: a description that persists long enough to be treated as a thing with a boundary.

If objects are induced, agents should be too. A companion paper on life \citep{six_birds_life} treats living systems as maintained, self-accounting closures rather than as a list of biological exceptions. Here we ask the parallel question for agency: what must be in place before something like an action exists, and once it exists, what does it mean for that action to be a cause?

\subsection{Agenthood and agency}
We separate two notions that are often run together.

\emph{Agenthood} is about enablement. Before there is an agent there are only micro-interactions. An agent exists when some package becomes stable enough to support an inside/outside distinction, a resource ledger, and an interface whose settings can be chosen. This is a claim about the existence and persistence of a level of description in which ``actions'' are well-defined variables.

\emph{Agency} is about causation within that level. Once the interface exists, we can ask whether different feasible settings of it lead to different outside futures. This is an interventional question, and we answer it with a capacity: how many distinguishable outside outcomes can be produced by choosing among affordable action sequences.

Keeping the two apart matters because many apparent signs of agency are modelling artefacts. A system with only one available action can move a great deal and still have no control. An external schedule can look like choice if it is wrongly typed as an action. We include explicit null tests for both mistakes.

\subsection{Thesis: an agent is a theory object}
In \SBT, the word \emph{theory} has a technical meaning: it is an induced level of description, specified by a lens (which distinctions can be made), a completion rule (an endomap whose fixed points are the objects the theory recognises), and an audit \citep{six_birds_theory}. It is not a hypothesis in the everyday sense. A \emph{theory object} is something that lives \emph{inside} such a level: a package that persists under the level's dynamics and accounting. For the finite controlled systems of this paper we specialise this to $T=(\Pi,L,\mathcal{F},B)$: a coarse-graining $\Pi$, induced dynamics $L$, feasibility rules $\mathcal{F}$, and resource bookkeeping $B$. This four-part form is our notation, motivated by \SBT\ rather than taken from it.

Our thesis is that an agent is a theory object of a particular kind: a maintained package whose interface supports a non-trivial, budget-respecting channel to outside futures. Agenthood is the existence and maintenance of the package; agency is the causal reach of its interface. In the finite setting of this paper, the controlled kernel together with its lenses, feasibility rule, and ledger plays the role of the theory, and the agent is the package that this theory makes stable.

\subsection{What we do}
We make the thesis testable with three computable measures on finite controlled Markov kernels (\cref{fig:overview}):
\begin{enumerate}[leftmargin=*, itemsep=0.25em]
  \item \textbf{Viability.} The viability kernel $\K$ is the largest set of safe states in which every state has an affordable action whose possible successors \emph{all} remain in the set. It is a greatest fixed point, and its size measures whether a package can be maintained at all.
  \item \textbf{Feasible empowerment.} For a horizon $H$, we form the channel from affordable action sequences to an outside variable at time $H$ and compute its capacity in bits. This measures difference-making.
  \item \textbf{Packaging defect.} For a coarse lens that hides some microscopic variables, we define a map $E$ that sends each coarse label to the label most likely to be observed $\tau$ steps later. The fraction of labels on which $E\circ E\neq E$ measures whether the coarse labels behave like stable objects.
\end{enumerate}
We apply these measures to a small ring world whose mechanisms can be switched on and off: a phase-dependent step (protocol), a paid repair action (maintenance), action costs and income (feasibility and accounting), and a skill variable that reduces slip (operator rewriting).

\begin{figure}[t]
\centering
% Overview diagram: the ring world (left) and the three measures (right).
\begin{tikzpicture}[
  font=\footnotesize,
  site/.style={circle, draw=black!55, fill=white, minimum size=4.4mm, inner sep=0pt},
  gain/.style={site, fill=black!15},
  agent/.style={circle, fill=blue!70!black, minimum size=2.6mm, inner sep=0pt},
  box/.style={draw=black!40, rounded corners=2pt, align=left, text width=5.4cm, inner sep=4pt},
  tag/.style={font=\scriptsize\bfseries, text=black!60, anchor=west},
  >={Stealth[length=1.8mm]}
]
% Ring world: sites numbered clockwise from the top; LEFT increases y.
\def\R{1.3}
\foreach \k in {0,...,7} {
  \pgfmathsetmacro{\ang}{90-45*\k}
  \ifnum\k=0
    \node[gain] (s\k) at (\ang:\R) {};
  \else
    \node[site] (s\k) at (\ang:\R) {};
  \fi
  \node[font=\scriptsize, text=black!55] at (\ang:\R+0.4) {$\k$};
}
\foreach \k [evaluate=\k as \j using {int(mod(\k+1,8))}] in {0,...,7} {
  \draw[black!30] (s\k) -- (s\j);
}
\node[agent] at (s2) {};
\draw[->, blue!70!black, thick] (-10:1.75) arc[start angle=-10, end angle=-38, radius=1.75];
\node[blue!70!black, font=\scriptsize, anchor=west] at (-26:1.8) {LEFT};
\draw[->, blue!70!black, thick] (10:1.75) arc[start angle=10, end angle=38, radius=1.75];
\node[blue!70!black, font=\scriptsize, anchor=west] at (26:1.8) {RIGHT};
\node[align=center, font=\scriptsize] at (0,0) {position $y$\\ phase $\phi$\\ damage $u$\\ ledger $r$};
\node[font=\scriptsize, text=black!55, align=center] at (0,-2.25) {shaded: income site};

% Measures
\begin{scope}[shift={(3.3,1.95)}]
  \node[tag] (ag) at (0,0) {AGENTHOOD: does a stable package exist?};
  \node[box, anchor=north west] (v) at (0,-0.25) {\textbf{Viability kernel} $\K$: largest safe set in which every state has an affordable action keeping all possible successors inside};
  \node[box, anchor=north west] (p) at ($(v.south west)+(0,-0.15)$) {\textbf{Packaging defect} $\Def(E)$: are the labels that $E$ produces fixed when $E$ is applied again?};
  \node[tag] (ac) at ($(p.south west)+(0,-0.3)$) {AGENCY: do its choices change the outside?};
  \node[box, anchor=north west] (e) at ($(ac.south west)+(0,-0.12)$) {\textbf{Feasible empowerment}: capacity from affordable action sequences to the position $y$ at horizon $H$};
\end{scope}
\end{tikzpicture}

\caption{Overview. \textbf{Left:} the ring world. An agent occupies one of $L=8$ sites. LEFT and RIGHT move it by one or two sites depending on a staged phase $\phi$ when the protocol is on; a hidden damage bit $u$ reverses the direction of motion; each action has a cost paid from a ledger $r$, which is replenished at income sites; REPAIR resets $u$ to $0$. \textbf{Right:} the three measures used in this paper. Viability and packaging ask whether a stable package exists (agenthood); feasible empowerment asks whether its interface choices change outside futures (agency).}
\label{fig:overview}
\end{figure}

\subsection{Main findings}
All results are exact or numerically converged computations on small, fully specified systems. They are existence witnesses under named configurations and policies, not general laws.
\begin{itemize}[leftmargin=*, itemsep=0.25em]
  \item \textbf{Null tests behave as they should.} A system with a single action has zero empowerment at every horizon tested. Treating an external schedule as an action manufactures a spurious 1~bit, which disappears under the correct model (\cref{sec:ex_nulls}).
  \item \textbf{Maintenance and objecthood.} In a funded ring world where repair is paid for every step, the packaging defect at a full phase cycle is 1 without repair and 0 with preventive repair, and the set of states that can be kept coherent grows from empty to 16. Two controls also reach zero defect without effective repair, so the defect alone does not detect maintenance; viability does (\cref{sec:ex_packaging}).
  \item \textbf{Order matters.} When the step size depends on the phase, one-step empowerment is unchanged (1.05~bits in both regimes), but the median over sampled viable states rises from 1.21 to 1.72~bits at two steps and from 1.29 to 1.85~bits at four and five steps (\cref{sec:ex_holonomy}).
  \item \textbf{Viability has a support threshold.} With coherence required for safety, any positive damage probability gives the same viability kernel, which shrinks with repair cost and vanishes once repair costs more than the available income. The transition is a change of support at zero noise plus a cost threshold, not a gradual decline (\cref{sec:ex_sweep}).
  \item \textbf{Skill increases difference-making.} Reducing slip through a skill variable raises median empowerment from 0.84 to 1.05 to 1.37~bits (\cref{sec:ex_learning}).
  \item \textbf{A mechanised anchor.} A Lean~4 development proves that the viability iteration, started from the safe set, stops after at most as many strict removals as there are safe states and returns the greatest safe controlled-invariant set (\cref{sec:repro,app:lean_viability}).
\end{itemize}

\subsection{Contributions}
\begin{itemize}[leftmargin=*, itemsep=0.25em]
  \item A precise reading of ``an agent is a theory object'' that separates agenthood (a maintained package) from agency (interface-to-outside difference-making), with each side tied to a computable measure.
  \item A small software pipeline for finite controlled kernels: robust viability kernels under ledger-gated feasibility, feasible empowerment computed with a Blahut--Arimoto iteration that stops only when explicit lower and upper capacity bounds agree to a set tolerance (in floating-point arithmetic), and a packaging-map diagnostic.
  \item A ring world with switchable mechanisms and a set of matched comparisons, together with null tests that guard against false positives.
  \item Reproducible artefacts: hashed configurations, a strict artefact auditor, exact rational cross-checks of key witnesses, and a Lean~4 proof of the finite fixed-point property used by the viability computation.
\end{itemize}

\subsection{Guide to the paper}
\Cref{sec:dictionary} maps the six mechanisms of \SBT\ to agency roles and to the measures we use. \Cref{sec:engine} defines the measures precisely and describes the ring world. \Cref{sec:ex_packaging,sec:ex_nulls,sec:ex_holonomy,sec:ex_ablations,sec:ex_sweep,sec:ex_learning} present the experiments. \Cref{sec:repro} explains how to reproduce and verify every number. \Cref{sec:discussion} discusses what the results do and do not show. \Cref{app:lean_viability} summarises the Lean development.
